\begin{Info}
% Universidade
{UNIVERSIDADE DO VALE DO ITAJAÍ}
% Escola
%{ESCOLA DO MAR, CIÊNCIA E TECNOLOGIA}
{ESCOLA POLITÉCNICA}
% Curso
{CURSO DE CIÊNCIA DA COMPUTAÇÃO}
% Titulo
{Análise da Otimização de CNNs 1D na Classificação de Imagens Hiperespectrais}
% Autor
{Eduardo Savian}
% Cidade e Data
{Itajaí (SC), Dezembro de 2026}
% Nome da Área de concentração
{Visão Computacional}
% Orientador(a)
{Felipe Viel, PhD.}
% Coorientador(a) <Nome do Coorientador(a)>, <Titulação> %%%%%%%% Se não tiver coorientador deixe vazio
{}
\end{Info}

%\begin{Dedicatoria}
% Dedicatória
% \end{Dedicatoria}

% \begin{Agradecimentos}
% Agradeço a todos.
% \end{Agradecimentos}

\begin{Epigrafe}
Todos os Homens, por natureza, desejam conhecer - Aristóteles
\end{Epigrafe}

\begin{Resumo}
A classificação de imagens hiperespectrais (HSI) com redes neurais convolucionais unidimensionais (CNNs 1D) é uma alternativa leve às arquiteturas que processam vizinhanças espaciais, e a fusão do identificador do superpixel às características espectrais compensa parte da falta de contexto espacial dessas redes. Este trabalho parte da CNN 1D com fusão do superpixel SEEDS de \citeonline{Souza2023} e a otimiza em quatro eixos (organização do código, algoritmos de segmentação, otimização de hiperparâmetros e paralelismo), medindo cada ganho de acurácia contra a configuração original, isto é, os hiperparâmetros de \citeonline{Souza2023} com a fusão SEEDS. O fluxo foi reimplementado em PyTorch como um pacote modular e determinístico, que reproduz o PCA, a divisão dos dados e os hiperparâmetros de treino do original e adota a CNN 1D do projeto, com quatro convoluções de 16 a 128 filtros, \textit{max-pooling} e densas de 256 e 128 neurônios; a parada antecipada segue a acurácia de validação, com teto de 200 épocas, e a CNN original ficou como comparação na linha de base. A avaliação usa Indian Pines, University of Pavia e Salinas Valley. O LSC e o Watershed foram calibrados para gerar o mesmo número de regiões que o SEEDS em cada cena. A busca em grade sobre o segmentador, a taxa de aprendizado e o \textit{dropout} foi comparada à otimização bayesiana com processo gaussiano, que percorre também o \textit{weight decay}, o lote, os componentes do PCA e a configuração do segmentador. No paralelismo, foram implementadas a partição da cena em faixas, com um modelo por partição em threads ou processos, e a fusão concorrente, em que uma thread segmenta enquanto um pool de threads extrai as características espectrais e faz a fusão. Com os hiperparâmetros do original e uma semente, a CNN do projeto obteve OA de 0,9142 em Indian Pines, contra 0,8960 da CNN original, e ficou 0,12 ponto percentual abaixo dela em University of Pavia e Salinas Valley. Com o mesmo número de regiões, o SEEDS gerou as regiões mais puras nas três cenas, e o Watershed segmentou de 28 a 62 vezes mais rápido. Em Indian Pines, com a CNN otimizada para cada segmentador e 5 sementes, o SEEDS (OA de 0,9324 $\pm$ 0,0062) superou o LSC (0,9152 $\pm$ 0,0045) por mais que o dobro do desvio-padrão combinado. \pendente{experimento 03: Watershed em Indian Pines; os três segmentadores em University of Pavia e Salinas Valley} \pendente{experimento 02: segmentadores com a configuração original fixa} \pendente{experimento 05: ganho da configuração otimizada sobre a original em Indian Pines, em OA, AA e Kappa} \pendente{experimento 04: grade contra bayesiana com 10, 40 e 200 épocas por \textit{trial}} \pendente{experimento 06: OA e tempo da partição contra o modelo único} \pendente{experimento 07: speedup do pool aquecido de processos sobre a execução sequencial} \pendente{experimento 08: speedup da fusão concorrente sobre o lote, por bloco e número de threads}

\textbf{Palavras-Chave}: Imagem Hiperespectral. Rede Neural Convolucional Unidimensional. Segmentação. Otimização de Hiperparâmetros. Processamento Paralelo.
\end{Resumo}

\begin{Abstract}
\textit{Hyperspectral image (HSI) classification with one-dimensional convolutional neural networks (1D CNNs) is a lightweight alternative to architectures that process spatial neighborhoods, and fusing the superpixel identifier with the spectral features compensates for part of the missing spatial context. This work starts from the 1D CNN with SEEDS superpixel fusion of Souza (2023) and optimizes it along four axes (code organization, segmentation algorithms, hyperparameter optimization, and parallelism), measuring each accuracy gain against the original configuration, that is, the hyperparameters of Souza (2023) with SEEDS fusion. The pipeline was reimplemented in PyTorch as a modular, deterministic package that reproduces the original PCA, data split, and training hyperparameters and adopts the project's 1D CNN, with four convolutional layers of 16 to 128 filters, max pooling, and dense layers of 256 and 128 neurons; early stopping follows validation accuracy with a ceiling of 200 epochs, and the original CNN is kept as a comparison in the baseline. The evaluation uses Indian Pines, University of Pavia, and Salinas Valley. LSC and Watershed were calibrated to produce the same number of regions as SEEDS on each scene. Grid search over the segmenter, learning rate, and dropout was compared with Bayesian optimization with a Gaussian process, which also covers weight decay, batch size, PCA components, and the segmenter configuration. For parallelism, two strategies were implemented: partitioning the scene into strips, with one model per partition in threads or processes, and concurrent fusion, in which one thread segments while a thread pool extracts the spectral features and performs the fusion. With the original hyperparameters and one seed, the project CNN reached an OA of 0.9142 on Indian Pines, against 0.8960 for the original CNN, and stayed 0.12 percentage points below it on University of Pavia and Salinas Valley. With the same number of regions, SEEDS produced the purest regions on all three scenes, and Watershed segmented 28 to 62 times faster. On Indian Pines, with the CNN tuned for each segmenter and 5 seeds, SEEDS (OA of 0.9324 $\pm$ 0.0062) outperformed LSC (0.9152 $\pm$ 0.0045) by more than twice the combined standard deviation.} \pendente{experimento 03: Watershed em Indian Pines; os três segmentadores em University of Pavia e Salinas Valley} \pendente{experimento 02: segmentadores com a configuração original fixa} \pendente{experimento 05: ganho da configuração otimizada sobre a original em Indian Pines, em OA, AA e Kappa} \pendente{experimento 04: grade contra bayesiana com 10, 40 e 200 épocas por \textit{trial}} \pendente{experimento 06: OA e tempo da partição contra o modelo único} \pendente{experimento 07: speedup do pool aquecido de processos sobre a execução sequencial} \pendente{experimento 08: speedup da fusão concorrente sobre o lote, por bloco e número de threads}

\textit{\textbf{Keywords}: Hyperspectral Image. One-Dimensional Convolutional Neural Network. Segmentation. Hyperparameter Optimization. Parallel Processing.}
\end{Abstract}

